\documentclass[11pt, letterpaper]{article}

% --- Essential Packages ---
\usepackage[utf8]{inputenc}
\usepackage[margin=1in]{geometry}
\usepackage{amsmath, amssymb, amsthm}
\usepackage{graphicx}
\usepackage{hyperref}

\begin{document}

% --- Custom Single Title Page ---
\begin{titlepage}
    \centering
    \vspace*{0.5in}

    \begin{figure}[htbp]
        \centering
        \includegraphics[width=0.3\textwidth]{"figures/University_of_Waterloo_seal.png"}
    \end{figure}

    \vspace{1.5cm}

    % Custom Title & Subtitle
    {\Huge \bfseries Lab 1: Modelling of the Motor System \par}

    \vspace{0.5cm}
    {\Large MTE 484: Digital Control Systems \par}
    \vspace{0.1cm}
    {\Large Instructor Sergio Livera \par}
    \vspace{0.5cm}
    {\large Department of Mechatronics and Mechanical Engineering \par}



    \vfill % Automatically fills vertical space, pushing author/date to the bottom

    % Custom Author Info
    {\Large \textbf{Levi Rogalla} \par}
    {\Large \textbf{Alex Veric} \par}

    \vspace{1.5cm}

    {\large \today \par}

    \vspace*{0.5in}
\end{titlepage}


\section{Introduction}
Lab one introduces the motor and gear mechanism that we will be using for the following
lab sessions in order to complete the ball and beam balancing project. The mechanism must
be properly characterized in order to facilitate the completion of this project. This
includes experimentally determining values for 'stiction' and the transfer function of the
motor, which will be described in detail in this lab report.

\section{Question C}
By this point in the lab, the set up was already complete. The Arduino was functioning as
intended, and with the sample code provided, it was possible to both send voltage values to
the motor and read sensor data, which included the motor rotational position and the ball's
position on the beam. These values were printed out on the serial monitor, allowing
them to be recorded.\par
The values of the motor position being output needed to be converted to radians in order
to be useful. Using the following formula, the
motor position output values could be linearized and converted into radians:
\begin{equation}
    p_{radians} = m * p_{sensor} + b
\end{equation}
Where the $m$ and $b$ values must be determined through a system of equations, using
experimental values of $p_{rad}$ and $p_{sensor}$ to do so. \par
The position of the motor was manually changed to three points: $0$, $\frac{\pi}{4}$,
and $-\frac{\pi}{4}$ radians. The values printed into the serial monitor were read at each
point, providing concrete values of $p_{rad}$ and $p_{sensor}$. These values were put
into the system of equations, and solving for this provided values for $m$ and $b$:
\begin{equation}
    m = -0.000359453251, b = 1.75789534
\end{equation}

\section{Question D}
The term 'stiction' describes the concept of static friction that exists in motors before
they start moving. This creates a range of voltage values that, when sent to the
motor, will not move the motor at all. This impacts the accuracy of the mechanism,
especially when considering the fine movements required to balance a ball on a beam.
This must be counteracted by creating a 'deadzone' at a certain range; for example,
if the deadzone range is determined to be -0.1 to 0.1 volts, any voltage input
within this range would be immediately rounded to either 0.1 or -0.1, depending on the
direction. \par
The values of voltage that determine the deadzone range must be experimentally determined.
Additionally, it is important to determine them in both directions, as the stiction in
motors is not often symmetrical. In order to complete this task, a function was derived
that, given a direction, methodically applies an increasing voltage to the motor, and
only stops the loop once the potentiometer detects the motor start to change position. \par


\end{document}